%-----------------------------------------------------------------------------
% Content for \input{work-in-progress}
% Paste everything below into work-in-progress.tex
%-----------------------------------------------------------------------------

\chapter{Work in Progress}

\newtheorem{rul}[subsection]{Rule}

%-----------------------------------------------------------------------------

%-----------------------------------------------------------------------------
\section{Circuit Simplification Rules}

%--- Rule 1 ------------------------------------------------------------------
\begin{rul}[Merge consecutive same-axis rotations]
\label{rul:merge-rot}
Two consecutive rotation gates of the same type on the same qubit compose into a
single gate whose angle is the sum of the individual angles:
\be
  R_P^{(q)}(\theta_1)\;R_P^{(q)}(\theta_2)
  \;=\;
  R_P^{(q)}(\theta_1+\theta_2),
  \qquad
  P \in \{X,Y,Z\}.
\ee
\end{rul}


\begin{rem}
Rotations about \emph{different} axes on the same qubit do not merge:
$R_Y(\theta_1)\,R_Z(\theta_2) \neq R_P(\theta)$ for any single Pauli $P$.
\end{rem}

%--- Rule 2 ------------------------------------------------------------------
\begin{rul}[Cancel consecutive identical CNOTs]
\label{rul:cnot-cancel}
Two consecutive CNOT gates with the same control $c$ and target $t$ cancel to the
identity:
\be
  \mathrm{CNOT}^{(c,t)}\;\mathrm{CNOT}^{(c,t)} \;=\; I.
\ee
\end{rul}

\begin{proof}
CNOT is its own inverse.  In the two-qubit computational basis ordered as
$|00\rangle,|01\rangle,|10\rangle,|11\rangle$ it reads
\be
  \mathrm{CNOT}
  =
  \begin{pmatrix}1&0&0&0\\0&1&0&0\\0&0&0&1\\0&0&1&0\end{pmatrix},
\ee
and a direct computation gives $\mathrm{CNOT}^2 = I_4$.
\end{proof}

%--- Rule 3 ------------------------------------------------------------------
\begin{rul}[Remove zero-angle rotations]
\label{rul:zero-angle}
A rotation gate whose angle satisfies $\theta \equiv 0 \pmod{4\pi}$ equals
$(\pm I)$ and may be dropped:
\be
  R_P^{(q)}(4\pi k) \;=\; (-1)^k\,I, \qquad k\in\Z.
\ee
In particular $R_P^{(q)}(0) = I$ and $R_P^{(q)}(4\pi) = I$.
The case $\theta = 2\pi$ gives $R_P(2\pi) = -I$, a global phase with no
observable consequence in unitary synthesis, and may also be removed.
\end{rul}

%--- Rule 4 ------------------------------------------------------------------
\begin{rul}[Angle reduction modulo $4\pi$]
\label{rul:angle-mod}
The rotation angle is $4\pi$-periodic:
\be
  R_P^{(q)}(\theta) \;=\; R_P^{(q)}(\theta \bmod 4\pi).
\ee
After applying Rule~\ref{rul:merge-rot}, the combined angle should be reduced to
the canonical interval $(-2\pi,\,2\pi]$ before testing Rule~\ref{rul:zero-angle}.
\end{rul}

%--- Rule 5 ------------------------------------------------------------------
\begin{rul}[SX periodicity]
\label{rul:sx-period}
Consecutive $\mathrm{SX}$ gates on the same qubit reduce as follows:
\begin{align}
  \mathrm{SX}^{(q)}\,\mathrm{SX}^{(q)}
    &\;=\; X^{(q)} \;=\; R_X^{(q)}(\pi), \\
  \bigl(\mathrm{SX}^{(q)}\bigr)^3
    &\;=\; \bigl(\mathrm{SX}^\dagger\bigr)^{(q)}, \\
  \bigl(\mathrm{SX}^{(q)}\bigr)^4
    &\;=\; I.
\end{align}
\end{rul}

\begin{proof}[Proof (two-gate case)]
The explicit matrix is
\be
  \mathrm{SX}
  =
  \frac{1}{2}\begin{pmatrix}1+i&1-i\\1-i&1+i\end{pmatrix},
  \qquad
  \mathrm{SX}^2
  =
  \begin{pmatrix}0&1\\1&0\end{pmatrix}
  = X. \qedhere
\ee
\end{proof}

\begin{rem}
A pair of consecutive $\mathrm{SX}$ gates is replaced by $R_X(\pi)$, which then
participates in further merging via Rule~\ref{rul:merge-rot}.  Four consecutive
$\mathrm{SX}$ gates are removed entirely.
\end{rem}

%--- Rule 6 ------------------------------------------------------------------
\begin{rul}[$R_Z$ commutes past the CNOT control]
\label{rul:rz-cnot}
A $Z$-rotation on the control qubit commutes with the CNOT:
\be
  R_Z^{(c)}(\theta)\;\mathrm{CNOT}^{(c,t)}
  \;=\;
  \mathrm{CNOT}^{(c,t)}\;R_Z^{(c)}(\theta).
\ee
\end{rul}


\begin{rem}
Rule~\ref{rul:rz-cnot} does not reduce gate count on its own.  Its role is to
move an $R_Z$ gate past a CNOT so that it becomes consecutive with another $R_Z$
on the same qubit, enabling Rule~\ref{rul:merge-rot} to fire.
\end{rem}

%--- Rule 7 ------------------------------------------------------------------
\begin{rul}[$R_X$ commutes past the CNOT target]
\label{rul:rx-cnot}
An $X$-rotation on the target qubit commutes with the CNOT:
\be
  R_X^{(t)}(\theta)\;\mathrm{CNOT}^{(c,t)}
  \;=\;
  \mathrm{CNOT}^{(c,t)}\;R_X^{(t)}(\theta).
\ee
\end{rul}


\begin{rem}
As with Rule~\ref{rul:rz-cnot}, this rule is used to expose consecutive $R_X$
pairs for merging rather than to directly reduce gate count.
\end{rem}

%-----------------------------------------------------------------------------
\section{Application Order}

The rules fall into three tiers.

\noindent\textbf{Tier 1 -- Structural cancellations} (gate count decreases directly):
Rule~\ref{rul:cnot-cancel} (CNOT pairs),
Rule~\ref{rul:zero-angle} (trivial rotations),
Rule~\ref{rul:sx-period} (SX runs).

\noindent\textbf{Tier 2 -- Merging} (two gates become one):
Rule~\ref{rul:merge-rot} combined with Rule~\ref{rul:angle-mod}.

\noindent\textbf{Tier 3 -- Commutation} (unlock further merges):
Rule~\ref{rul:rz-cnot} and Rule~\ref{rul:rx-cnot}.

A single-pass application order is:
\begin{enumerate}
  \item Scan left-to-right; apply Rule~\ref{rul:merge-rot} for every pair of
        consecutive same-type, same-qubit rotations.  Immediately follow with
        Rule~\ref{rul:angle-mod} and then Rule~\ref{rul:zero-angle}.
  \item Scan again; apply Rule~\ref{rul:cnot-cancel} for consecutive identical
        CNOTs.
  \item Scan again; apply Rule~\ref{rul:sx-period} for runs of consecutive
        $\mathrm{SX}$ gates.
  \item Optionally apply the commutation Rules~\ref{rul:rz-cnot}
        and~\ref{rul:rx-cnot} to expose new adjacent pairs, then repeat from
        step~1.
\end{enumerate}

Iterating until no rule fires gives a locally simplified circuit.  The result is
not guaranteed to be globally minimal, but removes the most common redundancies
introduced by the GQE token decoder.
